\scanpage{92}
\mspage{3}
% archive foliation: 91
de $\pi_1^t$ par $\overline I_0'$, soit
\deleted{$\unsure{\Pi'}$, \unsure{un}} $\Pi'$, et on a
\deleted{\unsure{donc}} \unsure{aussi} un \uline{épimorphisme}
\[
 \Pi'\longrightarrow\pi_1^t(\mathfrak X/D)
\]
prouver que c'est un \uline{isomorphisme}. ]

Ceci revient à vérifier ceci: étant donné un hom surjectif
\[
 \mu_0\longrightarrow J_0\longrightarrow0
\]
qui s'annule sur l'image de $\pi_2^t(U)$,
\deleted{\unsure{il}} $J_0$ apparaît comme groupe d'inertie dans
un revêtement \deleted{\unsure{tame}} modéré de $\mathfrak X/X_0$.
L'hypothèse signifie aussi que l'image de
$\beta\in H^2(\unsure{X|X_0},\mu)$ dans $H^2(\widetilde U^t,J_0)$
est nulle; elle est donc nulle déjà \deleted{\unsure{en}} à un
niveau fini [ N.B. \uline{par définition}, on définit la coh de
$\widetilde U^t$ comme une $\uline{\lim_{\longrightarrow}}$ \ldots,
\unsure{et} \deleted{$\unsure{U}$} \unsure{n'est pas noeth}
mais seulement loc. noeth \ldots ]. Donc quitte à faire
un revêtement modéré fini
\deleted{\unsure{de revêt} $\mathfrak X$ \unsure{nul} /
$D_i$ \unsure{disjoints} \unsure{on peut}}
\deleted{\unsure{loisir}}
\margin{\unsure{ce qui du} \unsure{pro.} \unsure{dans} est \unsure{faux}}
